% !TEX program = xelatex
\documentclass[12pt,a4paper]{article}

% ---------- Packages ----------
\usepackage[a4paper,margin=2.5cm]{geometry}
\usepackage{fontspec}
\usepackage{xeCJK}
\usepackage{setspace}
\usepackage{titlesec}
\usepackage{graphicx}
\usepackage{booktabs}
\usepackage{array}
\usepackage{tabularx}
\usepackage{longtable}
\usepackage{enumitem}
\usepackage{xcolor}
\usepackage{hyperref}
\usepackage{fancyhdr}
\usepackage{amsmath}
\usepackage{caption}
\usepackage{listings}
\usepackage{tocloft}
\usepackage{ragged2e}

% ---------- Font Settings ----------
\setmainfont{DejaVu Serif}
\setsansfont{DejaVu Sans}
\setmonofont{DejaVu Sans Mono}
\setCJKmainfont{AR PL UMing TW}
\setCJKsansfont{AR PL UMing TW}
\setCJKmonofont{AR PL UMing TW}
\XeTeXlinebreaklocale "zh"
\XeTeXlinebreakskip = 0pt plus 1pt

% ---------- Layout ----------
\onehalfspacing
\setlength{\parindent}{2em}
\setlength{\parskip}{0.25em}
\emergencystretch=3em
\sloppy
\renewcommand{\contentsname}{目錄}
\renewcommand{\refname}{參考文獻}
\renewcommand{\figurename}{圖}
\renewcommand{\tablename}{表}

\hypersetup{
  colorlinks=true,
  linkcolor=black,
  citecolor=black,
  urlcolor=blue,
  pdftitle={軟體供應鏈攻擊與建置來源證明},
  pdfauthor={NG, Chun Sing; OU, Chia-Yun}
}

\pagestyle{fancy}
\fancyhf{}
\lhead{軟體供應鏈攻擊與建置來源證明}
\rhead{\thepage}
\setlength{\headheight}{15pt}

\titleformat{\section}{\Large\bfseries}{\thesection}{1em}{}
\titleformat{\subsection}{\large\bfseries}{\thesubsection}{1em}{}
\titleformat{\subsubsection}{\normalsize\bfseries}{\thesubsubsection}{1em}{}

% ---------- Tables ----------
\newcolumntype{Y}{>{\RaggedRight\arraybackslash}X}
\newcolumntype{P}[1]{>{\RaggedRight\arraybackslash}p{#1}}

% ---------- Listings ----------
\lstdefinestyle{commandstyle}{
  basicstyle=\ttfamily\small,
  breaklines=true,
  frame=single,
  backgroundcolor=\color{gray!8},
  columns=fullflexible,
  keepspaces=true
}

% ---------- Title ----------
\title{\bfseries Software Supply Chain Attacks and Build Provenance\\
\large 軟體供應鏈攻擊與建置來源證明}
\author{NG, Chun Sing (r14922186) \and OU, Chia-Yun (r14922072)}
\date{\today}

\begin{document}

\maketitle

\begin{abstract}
現代軟體開發高度依賴開源套件、套件管理器、版本控制平台、持續整合與持續部署（CI/CD）流程，以及第三方套件註冊庫。這些機制使軟體重用與自動化發布更加便利，卻也使開發流程建立在複雜的信任關係之上。開發者在日常工作中執行 \texttt{npm install}、\texttt{pip install}、\texttt{mvn install} 或 \texttt{conda install} 時，實際上不只是下載程式碼，也可能間接信任大量 transitive dependencies、維護者帳號、註冊庫 metadata、安裝腳本、建置流程與發布管線。

本報告探討軟體供應鏈攻擊（software supply chain attacks）與建置來源證明（build provenance）在防禦此類攻擊時的角色。本文首先說明軟體供應鏈安全的定義、風險上升的原因，以及其在 OWASP Top 10:2025 中被列為 A03 Software Supply Chain Failures 的背景。接著，本文分析 Axios 與 TanStack 兩個 npm 生態系事件。Axios 事件顯示，攻擊者可透過被竊取的 maintainer credentials、惡意套件發布、隱藏相依套件與 \texttt{postinstall} 腳本，在一般安裝流程中植入惡意程式。TanStack 事件則進一步指出，即使 artifact 具有有效的 build provenance，只要受信任的 CI/CD workflow 或 cache trust boundary 遭到攻破，provenance 仍可能無法阻止惡意套件被合法流程發布。

本文的核心論點是：build provenance 是必要的，但不足以單獨構成完整防禦。Provenance 能夠提高 artifact traceability，使 source commit、builder identity、build workflow 與 artifact hash 之間的關係更可驗證；然而，它並不能證明 artifact 本身必然良性，也不能保證建置環境未被入侵。因此，有效的軟體供應鏈防禦應採取分層策略，結合 dependency scanning、lockfile discipline、install-time verification、maintainer account protection、CI/CD isolation、cache hardening、credential restriction、signed attestations 與 provenance verification。
\end{abstract}

\noindent\textbf{關鍵字：} software supply chain security, build provenance, SLSA, Sigstore, npm, CI/CD security, dependency attack

\newpage
\tableofcontents
\newpage

\section{Introduction}

現代軟體開發已不再是一個完全由單一開發團隊掌控的封閉流程。開發者經常透過套件管理器取得現成函式庫，並將建置、測試、發布與部署交給自動化工具執行。例如，日常開發中常見的指令如下：

\begin{lstlisting}[style=commandstyle]
npm install   pip install   mvn install   conda install
\end{lstlisting}

這些指令在工程上提高了效率，因為團隊可以快速重用已存在的軟體元件。然而，這種便利性同時擴大了攻擊面。安裝一個直接相依套件（direct dependency）時，套件管理器可能同時解析並下載數十個 transitive packages；在某些生態系中，安裝過程也可能執行 \texttt{preinstall}、\texttt{install} 或 \texttt{postinstall} 等 lifecycle scripts。換言之，開發者表面上只執行一個安裝指令，實際上卻信任了許多未經人工逐一檢查的上游元件與中介流程。

軟體供應鏈攻擊的危險性正在於此。攻擊者不一定需要直接攻破最終應用程式，而是可以攻擊已被開發者信任的任一環節，例如 dependency package、maintainer account、package registry、build cache、CI/CD workflow 或 release pipeline。一旦受信任的上游環節被攻破，惡意程式碼就可能透過正常安裝或正常發布流程進入大量 downstream projects。

本文聚焦的研究問題是：\textbf{build provenance 是否能有效防止軟體供應鏈攻擊？} 本報告採取較謹慎的立場。Build provenance 可以讓 artifact 的來源、建置流程與發布紀錄變得更可追蹤，也能提供 install-time verification 的基礎；但是，如果合法 workflow 本身已被攻擊者利用，provenance 仍可能產生一份「正確但不充分」的紀錄。因此，本報告將 provenance 視為 defense-in-depth 中的關鍵機制，而非單一完整解法。

\section{Background: Software Supply Chain Security}

\subsection{Definition of Software Supply Chain}

根據 NIST 的說明，software supply chain 是一組用來建立、轉換並評估軟體 artifacts 品質與政策符合性的步驟。在實務上，這條供應鏈涵蓋 source code、third-party dependencies、version control systems、build tools、CI/CD pipelines、package registries、artifact repositories、release workflows 與 deployment mechanisms。也就是說，現代軟體產品的安全性不只取決於開發團隊撰寫的程式碼，也取決於其所依賴的整個開發、建置與發布環境。

Software supply chain attack 則是指攻擊者攻破或操控供應鏈中的某個受信任節點。該節點可能是 dependency package、maintainer account、build cache、CI/CD workflow、registry account、signing process 或 release pipeline。攻擊者的目標不是單純利用應用程式本身的漏洞，而是濫用開發者、工具、套件與註冊庫之間既有的信任關係。由於熱門套件通常被大量專案重複使用，單一 upstream compromise 便可能造成大規模 downstream impact。

\subsection{Why Software Supply Chain Failure Became a Top Risk}

OWASP Top 10:2025 將 \textit{Software Supply Chain Failures} 排為 A03，顯示應用程式安全的風險重心正在轉變。傳統應用程式安全多半聚焦於應用程式內部缺陷，例如 injection、authentication failure、insecure design 或 access control failure。相較之下，software supply chain security 更關注應用程式是如何被建置、依賴哪些外部元件、由誰發布這些元件，以及最終 artifact 是否能被驗證。

OWASP Top 10:2025 並非完全由單一統計資料直接決定，而是 data-informed。其排名結合 contributed testing data、CVE-based exploit and impact scores，以及 community survey results。對於 A03:2025 Software Supply Chain Failures，OWASP 報告的數據包括 5.72\% average incidence rate、215,248 total occurrences，以及 11 個 mapped CVEs。這些數字應被解讀為目前測試資料與工具所能觀察到的風險表現，而不能直接等同於全球實際攻擊頻率。事實上，供應鏈攻擊往往涉及新發布惡意套件、憑證濫用或 CI/CD 流程被操控，這些情境未必會立即出現在 CVE 資料庫或傳統掃描結果中，因此真實風險可能被低估。

此問題之所以成為高風險議題，主要來自幾個結構性因素。第一，現代應用程式通常具有大型 dependency graph，少數 direct dependencies 可能帶入大量 transitive dependencies。第二，套件管理器高度自動化，會自動解析版本、下載套件、讀取 metadata，並在某些情況下執行安裝腳本。第三，CI/CD 已成為建置與發布的核心流程，使 workflow permissions、runner environment、OIDC tokens 與 build cache 成為重要攻擊面。第四，開發者通常不會人工審查每一個 package version、tarball、build log、artifact hash 與 provenance record。這些因素共同造成攻擊者得以透過一個受信任節點影響大量 downstream systems。

\subsection{Economic and Ecosystem Impact}

供應鏈攻擊的財務影響難以完全量化，因為許多事件不會公開揭露損失細節。然而，breach-cost data 仍顯示 supply chain compromise 的代價相當高。IBM 2025 breach-cost data 指出，third-party vendor and supply chain compromise 平均每次攻擊成本為 USD 4.91M，並被列為第二常見且第二高成本的 data breach vector。此類成本可能包括 incident response、downtime、credential rotation、CI/CD rebuild、legal risk 與 reputational damage。

同時，open-source package risk 也持續上升。ReversingLabs 2026 software supply chain security report 指出，2025 年 malicious open-source package detections 增加 73\%，其中 npm 約佔 detected OSS malware 的 90\%。這項趨勢說明，攻擊者已將 package ecosystems 視為具高度投資報酬率的攻擊面。npm 生態系尤其值得關注，因為它具有龐大的套件規模、快速的版本發布節奏、自動化安裝流程，以及高度 transitive 的 dependency structure。

\section{Threat Model and Attack Surfaces}

軟體供應鏈攻擊可以分為 dependency layer、source layer、build layer 與 release layer 四個層面觀察。每個層面具有不同的 attack surface，也需要不同類型的防禦控制。

\begin{table}[h]
\centering
\caption{軟體供應鏈中的主要攻擊面}
\begin{tabularx}{\textwidth}{P{3cm}Y}
\toprule
\textbf{Layer} & \textbf{Possible Attack Surface} \\
\midrule
Dependency & Malicious packages、typosquatting、dependency confusion、transitive dependency compromise，以及濫用 \texttt{preinstall}、\texttt{install}、\texttt{postinstall} lifecycle scripts。 \\
Source & Compromised maintainer accounts、malicious commits、review bypass、弱化的 protected branch settings，以及過度授權的 release permissions。 \\
Build & CI/CD workflow tampering、build cache poisoning、從 runner memory 擷取 token、不安全的 fork pull request workflows，以及 overprivileged OIDC tokens。 \\
Release & Registry account compromise、artifact substitution、missing signatures、missing attestations，以及缺少 provenance verification。 \\
\bottomrule
\end{tabularx}
\end{table}

在 npm 生態系中，lifecycle scripts 是特別重要的攻擊面。\texttt{preinstall}、\texttt{install} 與 \texttt{postinstall} 可能在套件安裝階段自動執行。因此，攻擊者若能發布含有惡意 lifecycle script 的套件版本，受害者可能在 application import 該套件之前就已經執行惡意程式碼。這使 installation phase 本身成為供應鏈攻擊的重要 execution vector。

\section{Case Study I: Axios npm Compromise}

\subsection{Incident Overview}

Axios 是 JavaScript 生態系中廣泛使用的 HTTP client。根據簡報資料，Axios 每週約有 100 million downloads。2026 年 3 月 31 日，Axios 發生 npm supply chain compromise。攻擊者使用被竊取的 maintainer credentials 發布惡意套件版本，暴露時間約為三小時。惡意版本注入一個 attacker-controlled hidden dependency，並透過 \texttt{postinstall} script 在正常安裝過程中安裝 Remote Access Trojan（RAT）。

此事件具有代表性，因為攻擊並不依賴使用者主動執行應用程式。使用者只要執行 \texttt{npm install axios}，安裝流程就可能觸發惡意腳本。換言之，攻擊發生於開發與安裝階段，而不是部署後的 runtime 階段。這種攻擊方式也使傳統 application-level security controls 較難觀察到最初的 compromise path。

根據簡報內容，此事件也影響 OpenAI 的 macOS app-signing pipeline，使 ChatGPT Desktop 與 Codex 的 code-signing certificates 面臨風險。這一點凸顯供應鏈攻擊的放大效應：當受影響者不是單一終端使用者，而是擁有簽章憑證、建置權限或發布權限的開發環境時，攻擊後果可能進一步擴散到更多 downstream users。

\subsection{Attack Flow}

Axios 事件可概念化為下列攻擊流程：

\begin{lstlisting}[style=commandstyle]
npm install axios
  -> malicious axios version
      -> attacker-controlled hidden dependency
          -> postinstall script executes
              -> Remote Access Trojan installed silently
\end{lstlisting}

此流程的危險性在於它與正常開發行為高度相似。套件名稱、registry entry 與安裝指令都可能看起來正常，但實際 tarball 與 dependency graph 已被攻擊者操控。由於沒有明顯 prompt 或 manual review，受害者很可能在毫無察覺的情況下被 compromise。

\subsection{Limitations of Common Defenses}

Axios 事件也揭示常見 dependency defenses 的限制。\texttt{npm audit} 與 Snyk 對已知 CVEs 很有價值，但新發布的 malicious version 可能尚未被登錄於 vulnerability database。\texttt{npm ci} 與 lockfile 能夠改善 reproducibility，並驗證 registry 中的 pinned versions 與 hash；然而，如果 lockfile 是在惡意版本可用期間產生，它可能反而固定 compromised version。Dependabot 則能自動更新已知 vulnerable dependencies，但其判斷仍高度依賴 CVE 與 vulnerability feeds；自動升級也可能意外引入剛被攻破的新版本。

\begin{longtable}{P{3.3cm}P{5.2cm}P{6.1cm}}
\caption{常見防禦方式與其限制}\\
\toprule
\textbf{Defense} & \textbf{Protection} & \textbf{Limitation} \\
\midrule
\endfirsthead
\toprule
\textbf{Defense} & \textbf{Protection} & \textbf{Limitation} \\
\midrule
\endhead
\texttt{npm audit} / Snyk & 偵測 dependency tree 中已知 CVEs 與已知 vulnerable dependencies。 & 對 newly published malicious versions 反應較慢，因為惡意版本可能尚未被分配 CVE 或收錄進資料庫。 \\
\addlinespace
\texttt{npm ci} + lockfile & 對 pinned versions 與 registry hash 進行驗證，提高 build reproducibility。 & 仍然信任 lockfile；若 lockfile 於攻擊暴露期間產生，可能鎖定惡意版本。 \\
\addlinespace
Dependabot & 自動更新已知 vulnerable dependencies。 & 仍依賴 vulnerability database；更新本身也可能引入新風險。 \\
\bottomrule
\end{longtable}

因此，供應鏈防禦不能只依靠 known-vulnerability scanning。對於惡意套件發布、憑證竊取與安裝腳本濫用，防禦也必須納入 package metadata、publishing behavior、lifecycle scripts、provenance status 與 install-time risk checks。

\section{The Core Gap: No Tamper-Evident Chain}

Axios 事件反映一個更根本的問題：在 source code、build system、registry package 與 final installation 之間，傳統套件安裝流程未必形成可驗證的 tamper-evident chain。一般開發者在安裝套件時，通常無法清楚回答以下問題：

\begin{itemize}
  \item Artifact 是否由預期的 source commit 建置而來？
  \item Build steps 是否未被修改或濫用？
  \item Registry 中發布的 tarball 是否與 CI output 相符？
  \item Artifact 是否由預期的 workflow、builder identity 與 release process 產生？
\end{itemize}

當這些問題無法被回答時，攻擊可能維持沉默。套件的名稱、版本號與 registry 頁面仍然可信，但 artifact 本身未必來自預期來源。Build provenance 的主要價值，就是試圖在 source 與 install 之間建立可簽署、可驗證、可追蹤的證據鏈。

\section{Build Provenance as a Defense Mechanism}

\subsection{Concept of Build Provenance}

Build provenance 是描述 software artifact 如何被產生的 metadata。它通常記錄 source repository、commit hash、builder identity、build workflow、build time、build type 與 artifact hash。簡報中以一句話概括其意義：\textit{Built from commit X, by pipeline Y, at time T}。若這份紀錄經由簽署並寫入 transparency log，consumer 便能在安裝或驗證時檢查 artifact 是否符合預期。

在 SLSA 與 Sigstore 的模型中，build provenance 通常結合 OIDC identity、keyless signing、artifact hashes 與 Rekor public append-only log。SLSA 提供供應鏈完整性框架，說明 provenance 應證明什麼；Sigstore 則提供簽署與 transparency log infrastructure，使 provenance statement 更具可驗證性。這些機制的共同目標不是單純產生 metadata，而是使 source-to-artifact path 更具 tamper-evident 特性。

\subsection{What Provenance Can Prove}

Build provenance 可以協助回答三類核心問題。第一，它可透過 signed commit hash 記錄 artifact 是否來自預期 source commit。第二，它可透過 OIDC-bound builder identity 連結 artifact 與特定 build pipeline。第三，它可透過 artifact hash 檢查 registry tarball 是否與 attested CI output 一致。

\begin{table}[h]
\centering
\caption{Build provenance 可以提供的驗證資訊}
\begin{tabularx}{\textwidth}{P{4.7cm}Y}
\toprule
\textbf{Question} & \textbf{How Provenance Helps} \\
\midrule
Correct source commit? & Commit hash 會被記錄並簽署，使 consumer 能檢查 artifact 是否來自預期 source revision。 \\
Build steps unmodified? & Pipeline identity 可透過 OIDC 與 attestation 綁定，使 consumer 能檢查 builder 是否符合預期。 \\
Tarball matches CI output? & Artifact hash 可被記錄在 provenance statement 或 transparency log 中，使 registry artifact 可與 CI output 比對。 \\
\bottomrule
\end{tabularx}
\end{table}

因此，provenance 的核心貢獻在於 traceability。Consumer 不必只依賴 package name 與 version number，而是可以要求 artifact 提供來源、建置流程與簽章證據。

\subsection{Verification: Turning the Record into Prevention}

Provenance 本身是一份紀錄；若要將其轉化為預防性控制，必須在 install time 或 policy enforcement 階段進行驗證。SLSA 的 artifact verification 模型可簡化為三個步驟。第一，consumer 驗證 provenance signature，確認 attestation 未被竄改。第二，consumer 將 provenance 中的 source repository、builder identity 與 build type 與 expected values 比較。第三，在需要更高保證的情境下，consumer 可遞迴驗證 dependencies。

Expected values 可以由 maintainer 在 repository 中宣告，也可以由 organization policy 指定，或由 first observed attestation 學習而來。這個設計使 provenance 不只是被動紀錄，而是可以成為 install-time gate。例如，如果某 package 先前版本沒有 attestation，工具可以發出 warning；如果先前版本有 attestation，但目標版本突然缺少 attestation，工具可以 raise error 或阻止安裝；若目標版本有 attestation，工具則可以驗證 signature 並繼續安裝流程。

\subsection{Technology Stack and Workflow}

實務上，build provenance 涉及多個技術層次。SLSA 定義 provenance 與 build levels 需要證明的內容；Sigstore 提供 keyless signing 與 Rekor transparency log；registry 與 CI platform 負責產生、儲存與發布 attestations；consumer-side tools 如 npq 或 Socket Firewall 則可在安裝前進行 package risk checks 與 provenance verification。

\begin{table}[h]
\centering
\caption{Build provenance 相關技術層次}
\begin{tabularx}{\textwidth}{P{3.2cm}P{4cm}Y}
\toprule
\textbf{Layer} & \textbf{Technology} & \textbf{Role} \\
\midrule
Framework & SLSA & 定義 provenance 應證明的內容，以及 build integrity levels。 \\
Cryptography & Sigstore & 提供 keyless signing 與 public transparency log support，包括 Rekor。 \\
Platform & Registry + CI & 產生、儲存與發布 attestations，例如 npm、PyPI 或 GitHub Actions。 \\
Consumer & npq / Socket Firewall & 作為 pre-install gate，檢查 package risk 與 provenance status。 \\
\bottomrule
\end{tabularx}
\end{table}

Producer side 的流程通常是在 CI/CD 中從特定 source commit 建置 artifact，產生 provenance statement，將 commit、builder、artifact hash 等資訊簽署並記錄，最後由 registry 儲存 package 與 attestation metadata。Consumer side 則是在安裝前取得 package 與 attestations，驗證 signature，檢查 expected source、builder 與 build type，並依政策決定 warning、blocking 或允許安裝。

\subsection{Adoption Status}

Provenance support 已逐漸出現在主要 package ecosystems 中，但採用仍屬早期階段。簡報資料整理指出，npm 的 \texttt{npm publish --provenance} 自 2023 年 10 月起 generally available；GitHub Artifact Attestations 自 2024 年 6 月起 generally available；PyPI 於 2024 年 11 月支援 digital attestations；Maven、RubyGems 與 Conda 也陸續開始支援或實驗相關機制。

\begin{longtable}{P{2.7cm}P{4.3cm}P{4cm}P{4.2cm}}
\caption{Provenance-related mechanisms 的採用狀態}\\
\toprule
\textbf{Ecosystem} & \textbf{Status} & \textbf{Adoption} & \textbf{Source Mentioned in Slides} \\
\midrule
\endfirsthead
\toprule
\textbf{Ecosystem} & \textbf{Status} & \textbf{Adoption} & \textbf{Source Mentioned in Slides} \\
\midrule
\endhead
npm & 2023 年 10 月起 generally available。 & beta 期間已有 3,800 packages；provenance-enabled packages 超過 500M downloads。 & GitHub Blog / Sigstore Blog \\
\addlinespace
PyPI & 2024 年 11 月起 generally available。 & 到 2025 年底，17\% uploads included attestations。 & PyPI Blog \\
\addlinespace
GitHub Attestations & 2024 年 6 月起 generally available。 & 簡報未列出公開 adoption figure。 & GitHub Blog \\
\addlinespace
Maven & 2025 年 1 月起 opt-in。 & Early adopters only。 & Sonatype \\
\addlinespace
RubyGems & In progress。 & 2025 年 3 月有 20 top gems shipping attestations。 & RubyGems Blog \\
\addlinespace
Conda & 2025 年 8 月起 beta。 & 簡報未列出公開 adoption figure。 & prefix.dev \\
\bottomrule
\end{longtable}

此採用狀態說明 provenance 正在逐漸成熟，但尚未形成普遍強制標準。主要限制不只在於工具是否支援，也包括 maintainer awareness、workflow migration overhead，以及 consumer 是否真的在 install time 驗證 provenance，而不是僅僅看到 attestation 存在就假設安全。

\section{Case Study II: TanStack npm Supply-Chain Compromise}

\subsection{Incident Overview}

TanStack 事件凸顯 build provenance 的邊界。根據簡報資料，TanStack 於 2026 年 5 月 11 日發生 npm supply-chain compromise。短時間內，42 packages 與 84 malicious versions 被發布；payload 目標包括 cloud credentials、API tokens 與 SSH keys。此事件由外部研究者在 26 分鐘內偵測到。

此事件最值得注意之處在於：惡意版本具有 valid build provenance。也就是說，問題不在於 provenance 沒有產生，也不在於 attestation 與 package 完全不相符。相反地，provenance 正確記錄 package 由 TanStack 的合法 workflow 產生；真正的問題是合法 workflow 所依賴的 cache 與 token trust boundary 已被攻擊者濫用。

\subsection{How the Attack Worked}

簡報中的攻擊流程可概括如下：攻擊者透過 fork pull request 毒化 shared CI cache。當後續合法 release workflow 執行時，它 restore 了被毒化的 cache。被 restore 的 cache 接著從 runner memory 中擷取 OIDC token，並透過合法發布流程推出 malicious packages。簡化後的流程如下：

\begin{lstlisting}[style=commandstyle]
fork pull request
  -> shared CI cache poisoned
      -> legitimate release workflow restores cache
          -> OIDC token extracted from runner memory
              -> malicious packages published through legitimate workflow
                  -> valid build provenance generated
\end{lstlisting}

此事件的關鍵在於 trust boundary 的錯置。Fork pull request workflow 屬於較低信任等級，而 release workflow 屬於高信任等級；若兩者共享 cache 或可互相影響，攻擊者便可能利用低信任入口污染高信任發布流程。這說明 CI/CD cache、runner memory、OIDC token 與 workflow permissions 都是 provenance 之外必須保護的攻擊面。

\subsection{Why This Challenges Provenance}

TanStack 事件證明了一個重要限制：\textbf{provenance proves origin, not innocence}。Provenance 可以證明 artifact 來自某個 repository、workflow 或 builder，但它無法單獨證明 workflow 執行時沒有被污染。若合法 workflow 本身被攻擊者利用，provenance 仍會準確記錄該 artifact 來自該 workflow，但 artifact 依然可能是惡意的。

因此，build provenance 是必要條件，而不是充分條件。它能使供應鏈更透明、更可追蹤，也能幫助 consumer 發現 artifact substitution 或 missing attestation。然而，它不能取代 CI/CD isolation、cache hardening、credential restriction 與 workflow permission auditing。換言之，provenance 提高了攻擊成本，但攻擊者仍可能透過破壞 build environment 來繞過 provenance 所提供的保證。

\subsection{Industry Response: OpenAI}

簡報也提到 OpenAI 對 TanStack 事件的回應。TanStack incident 中，有兩台 OpenAI employee devices 受到影響，而當時相關控制措施仍處於 phased rollout。Axios 事件後已部署的控制措施包括：強化 CI/CD pipelines 中的 credential materials、設定 package manager 的 \texttt{minimumReleaseAge} 以拒絕剛發布的新版本，以及使用 security software validate provenance of new packages。

此回應與前述分析一致。Provenance verification 很重要，但組織仍需要降低「剛發布惡意版本」進入開發環境的機率，也需要保護 CI/CD 與開發設備中的 credentials。這些控制措施共同構成 layered defense，而非依賴單一機制解決所有風險。

\section{Toward a Layered Defense Strategy}

Axios 與 TanStack 兩個案例分別呈現不同層面的供應鏈風險。Axios 主要說明 maintainer credentials、malicious package publication 與 lifecycle scripts 的問題；TanStack 則說明 CI/CD cache、OIDC tokens 與 trusted workflow compromise 的問題。兩者共同指出：單一防禦機制無法覆蓋完整供應鏈攻擊面。因此，實務防禦必須採取分層策略。

\subsection{Dependency Layer}

Dependency layer 的目標是降低惡意或脆弱套件進入專案的機率。實務上，組織應使用 lockfiles 控制版本與 hash，並避免不必要的 dependency updates；同時，可使用 SBOMs 記錄 dependency composition，並整合 \texttt{npm audit}、OSV 或 Snyk 等 vulnerability scanners 以偵測已知 vulnerabilities。然而，這些工具不足以處理 newly published malicious packages，因此也應搭配 npq 或 Socket Firewall 等 install-time gates，檢查 package metadata、provenance status 與 suspicious behavior。對於高風險環境，minimum release age policy 也能降低剛發布惡意版本立即進入專案的機率。

\subsection{Source Layer}

Source layer 主要保護 source repositories 與 maintainer accounts。維護者帳號若被竊取，攻擊者可直接發布惡意版本或修改 release process。因此，multi-factor authentication、protected branches、mandatory code review 與 restricted publishing permissions 都是必要控制。除此之外，組織也應監控異常 commits、releases、token usage 與 maintainer activity，以便在受信任帳號行為異常時及早發現風險。

\subsection{Build Layer}

Build layer 是 TanStack 事件最重要的教訓。組織必須嚴格分離 untrusted pull request workflows 與 trusted release workflows，避免 fork PR workflows 與 release workflows 共用高信任 cache。GitHub Actions tokens、OIDC tokens 與 registry tokens 應遵循 least privilege principle，並避免 credential materials 出現在 runner memory、logs 或可被低信任流程讀取的位置。Cache restore paths、build scripts 與 release scripts 也應被視為安全邊界的一部分，而不是單純的效能最佳化工具。

\subsection{Release Layer}

Release layer 的目標是在 publication 與 installation 階段使 artifacts 可驗證。組織應使用 signed attestations 與 build provenance，將 source commit、builder identity、artifact hash 與 build workflow 記錄於可驗證 metadata 中。更重要的是，consumer 應在 install time 驗證 provenance，而不是只在 publish time 產生 provenance。若某 package 先前具有 attestation，但新版本突然缺少 attestation，安裝工具應至少 warning，甚至依政策 block installation。

\section{Discussion}

\subsection{Practical Value of Build Provenance}

Build provenance 的實務價值在於將原本隱性的信任關係轉化為可驗證的 evidence。過去，consumer 往往只能看到 package name、version number 與 registry metadata，卻難以確認 artifact 是否真的來自預期 source code 與 build process。Provenance 透過 signed attestations 建立 source、builder 與 artifact 之間的關係，使 artifact substitution、unauthorized publishing，以及部分 maintainer-compromise scenarios 更容易被偵測。

此外，provenance 也為 policy enforcement 提供基礎。企業或組織可以要求安裝的 packages 必須來自特定 repository、特定 builder 或具有有效 attestation。這使供應鏈安全從「相信 package 應該是正確的」轉為「要求 package 提供可機器驗證的證據」。從這個角度看，provenance 不只是紀錄工具，也是一種可被自動化安全政策利用的基礎建設。

\subsection{Fundamental Limitation of Build Provenance}

然而，provenance 不應被誤解為 malware scanner，也不應被視為完整的 build environment integrity proof。TanStack 事件說明，若合法 workflow 被污染，provenance 可能仍然完全有效，但其所證明的是 artifact 的來源，而不是 artifact 的良性。換言之，provenance 可以回答「這個 artifact 從哪裡來」，但不能單獨回答「這個 artifact 是否安全」。

因此，較精確的結論是：\textbf{build provenance raises the bar, but the bar can still be cleared}。它提高攻擊者偽造 artifact origin 的難度，也使 missing provenance 或 unexpected builder 更容易被發現；但攻擊者仍可能透過 CI/CD cache poisoning、credential theft、workflow abuse 或 maintainer account compromise 繞過部分保證。Provenance 應被採用並驗證，但必須與 CI/CD isolation、cache hardening、credential restriction、dependency scanning 與 install-time controls 一起使用。

\subsection{Synthesis of the Two Case Studies}

Axios 與 TanStack 的差異有助於理解供應鏈安全的完整範圍。Axios 事件偏向 package publication 與 lifecycle script abuse：攻擊者利用被竊取的 maintainer credentials 發布惡意套件版本，並透過安裝腳本完成入侵。此情境中，provenance 與 install-time verification 可以增加使用者在安裝前發現異常的機會。

TanStack 事件則偏向 build pipeline compromise：攻擊者不是單純偽造 artifact origin，而是讓合法 workflow 產生惡意 artifact。此情境中，provenance 仍會產生看似正確的 attestation。因此，TanStack 事件提醒我們，build provenance 的可信度取決於 build environment 的可信度。若 build environment 的 trust boundary 設計不良，provenance 只能忠實記錄一個已被污染的流程。

\section{Conclusion}

Software supply chain attacks 之所以危險，是因為它們利用了現代軟體開發效率背後的信任結構。開發者信任套件管理器、套件維護者、註冊庫、CI/CD workflows、build caches 與 automated release processes。攻擊者則選擇攻擊這些受信任節點，使惡意程式碼能以正常開發活動的外觀進入 downstream systems。

Axios 事件顯示，被竊取的 maintainer credentials、malicious package publication、hidden dependencies 與 \texttt{postinstall} scripts 可以在正常安裝流程中完成入侵。TanStack 事件則進一步顯示，即使惡意套件具有 valid build provenance，只要 CI/CD workflow 或 shared cache 已被攻破，provenance 仍不足以阻止攻擊。兩個案例共同支持本文主張：build provenance 是必要的，但不是充分的。

Build provenance 的真正價值在於使 source、builder 與 artifact 之間的關係更可見、更可簽署、更可驗證。它可以支援 install-time verification 與 policy enforcement，並提高 artifact substitution 或 missing attestation 的可見度。然而，它不能取代 secure build environment，也不能單獨證明 artifact 是良性的。因此，未來的軟體供應鏈防禦應將 provenance 視為基礎機制，並同時部署 dependency scanning、lockfile discipline、minimum release age policy、maintainer account protection、CI/CD isolation、cache hardening、credential restriction 與 provenance verification。只有透過這種 layered defense，才能較合理地降低現代軟體供應鏈攻擊的系統性風險。

\newpage
\begin{thebibliography}{99}

\bibitem{nist204d}
National Institute of Standards and Technology.
\textit{NIST SP 800-204D: Strategies for the Integration of Software Supply Chain Security in DevSecOps CI/CD Pipelines}.
\url{https://nvlpubs.nist.gov/nistpubs/SpecialPublications/NIST.SP.800-204D.pdf}

\bibitem{owasp-cheat-sheet}
OWASP.
\textit{Software Supply Chain Security Cheat Sheet}.
\url{https://cheatsheetseries.owasp.org/cheatsheets/Software_Supply_Chain_Security_Cheat_Sheet.html}

\bibitem{owasp-top10-a03}
OWASP.
\textit{A03:2025 Software Supply Chain Failures}.
\url{https://owasp.org/Top10/2025/A03_2025-Software_Supply_Chain_Failures/}

\bibitem{owasp-methodology}
OWASP.
\textit{OWASP Top 10:2025 Introduction and Methodology}.
\url{https://owasp.org/Top10/2025/0x00_2025-Introduction/}

\bibitem{slsa-threats}
SLSA.
\textit{Supply Chain Threats}.
\url{https://slsa.dev/spec/v1.0/threats-overview}

\bibitem{slsa-about}
SLSA.
\textit{About SLSA}.
\url{https://slsa.dev/spec/v1.2/about}

\bibitem{slsa-verify}
SLSA.
\textit{Verifying Artifacts}.
\url{https://slsa.dev/spec/v1.0/verifying-artifacts}

\bibitem{sigstore}
OpenSSF.
\textit{Sigstore}.
\url{https://www.sigstore.dev/}

\bibitem{npq}
Liran Tal.
\textit{npq: Safely install packages with npm or yarn by auditing them as part of your install process}.
\url{https://github.com/lirantal/npq}

\bibitem{socket-firewall}
Socket.
\textit{Introducing Socket Firewall}.
\url{https://socket.dev/blog/introducing-socket-firewall}

\bibitem{axios-postmortem}
Axios.
\textit{Post Mortem: axios npm supply chain compromise}.
\url{https://github.com/axios/axios/issues/10636}

\bibitem{microsoft-axios}
Microsoft Security Blog.
\textit{Mitigating the Axios npm supply chain compromise}.
\url{https://www.microsoft.com/en-us/security/blog/2026/04/01/mitigating-the-axios-npm-supply-chain-compromise/}

\bibitem{openai-axios}
OpenAI.
\textit{Our response to the Axios developer tool compromise}.
\url{https://openai.com/index/axios-developer-tool-compromise}

\bibitem{tanstack-postmortem}
TanStack.
\textit{Postmortem: TanStack npm supply-chain compromise}.
\url{https://tanstack.com/blog/npm-supply-chain-compromise-postmortem}

\bibitem{tanstack-followup}
TanStack.
\textit{Hardening TanStack After the npm Compromise}.
\url{https://tanstack.com/blog/incident-followup}

\bibitem{openai-tanstack}
OpenAI.
\textit{Our response to the TanStack npm supply chain attack}.
\url{https://openai.com/index/our-response-to-the-tanstack-npm-supply-chain-attack}

\bibitem{github-attestations}
GitHub Blog.
\textit{Artifact Attestations is generally available}.
\url{https://github.blog/changelog/2024-06-25-artifact-attestations-is-generally-available/}

\bibitem{npm-provenance}
GitHub Blog.
\textit{Introducing npm package provenance}.
\url{https://github.blog/security/supply-chain-security/introducing-npm-package-provenance/}

\bibitem{sigstore-npm-ga}
Sigstore Blog.
\textit{npm's Sigstore-powered provenance goes GA}.
\url{https://blog.sigstore.dev/npm-provenance-ga/}

\bibitem{pypi-attestations}
PyPI Blog.
\textit{PyPI now supports digital attestations}.
\url{https://blog.pypi.org/posts/2024-11-14-pypi-now-supports-digital-attestations/}

\bibitem{pypi-2025}
PyPI Blog.
\textit{PyPI 2025 Year in Review}.
\url{https://blog.pypi.org/posts/2025-12-31-pypi-2025-in-review/}

\bibitem{sonatype-sigstore}
Sonatype Central.
\textit{Sigstore Signature Validation via Portal}.
\url{https://central.sonatype.org/news/20250128_sigstore_signature_validation_via_portal/}

\bibitem{rubygems-updates}
RubyGems Blog.
\textit{February 2025 Updates}.
\url{https://blog.rubygems.org/2025/03/19/february-rubygems-updates.html}

\bibitem{conda-sigstore}
prefix.dev.
\textit{Securing the Conda package supply chain with Sigstore}.
\url{https://prefix.dev/blog/securing-the-conda-package-supply-chain-with-sigstore}

\bibitem{ibm-cost}
IBM.
\textit{Cost of a Data Breach Report 2025 / Attack vector overview}.
\url{https://www.ibm.com/think/topics/attack-vector}

\bibitem{reversinglabs}
ReversingLabs.
\textit{Software Supply Chain Security Report 2026}.
\url{https://www.reversinglabs.com/resources/software-supply-chain-security-report-2026}

\end{thebibliography}

\end{document}
